% Options for packages loaded elsewhere
\PassOptionsToPackage{unicode}{hyperref}
\PassOptionsToPackage{hyphens}{url}
\documentclass[
]{article}
\usepackage{xcolor}
\usepackage[margin=1in]{geometry}
\usepackage{amsmath,amssymb}
\setcounter{secnumdepth}{-\maxdimen} % remove section numbering
\usepackage{iftex}
\ifPDFTeX
  \usepackage[T1]{fontenc}
  \usepackage[utf8]{inputenc}
  \usepackage{textcomp} % provide euro and other symbols
\else % if luatex or xetex
  \usepackage{unicode-math} % this also loads fontspec
  \defaultfontfeatures{Scale=MatchLowercase}
  \defaultfontfeatures[\rmfamily]{Ligatures=TeX,Scale=1}
\fi
\usepackage{lmodern}
\ifPDFTeX\else
  % xetex/luatex font selection
\fi
% Use upquote if available, for straight quotes in verbatim environments
\IfFileExists{upquote.sty}{\usepackage{upquote}}{}
\IfFileExists{microtype.sty}{% use microtype if available
  \usepackage[]{microtype}
  \UseMicrotypeSet[protrusion]{basicmath} % disable protrusion for tt fonts
}{}
\makeatletter
\@ifundefined{KOMAClassName}{% if non-KOMA class
  \IfFileExists{parskip.sty}{%
    \usepackage{parskip}
  }{% else
    \setlength{\parindent}{0pt}
    \setlength{\parskip}{6pt plus 2pt minus 1pt}}
}{% if KOMA class
  \KOMAoptions{parskip=half}}
\makeatother
\usepackage{graphicx}
\makeatletter
\newsavebox\pandoc@box
\newcommand*\pandocbounded[1]{% scales image to fit in text height/width
  \sbox\pandoc@box{#1}%
  \Gscale@div\@tempa{\textheight}{\dimexpr\ht\pandoc@box+\dp\pandoc@box\relax}%
  \Gscale@div\@tempb{\linewidth}{\wd\pandoc@box}%
  \ifdim\@tempb\p@<\@tempa\p@\let\@tempa\@tempb\fi% select the smaller of both
  \ifdim\@tempa\p@<\p@\scalebox{\@tempa}{\usebox\pandoc@box}%
  \else\usebox{\pandoc@box}%
  \fi%
}
% Set default figure placement to htbp
\def\fps@figure{htbp}
\makeatother
\setlength{\emergencystretch}{3em} % prevent overfull lines
\providecommand{\tightlist}{%
  \setlength{\itemsep}{0pt}\setlength{\parskip}{0pt}}
\usepackage{float}
\floatplacement{figure}{H}
\usepackage{booktabs}
\usepackage{longtable}
\usepackage{array}
\usepackage{multirow}
\usepackage{wrapfig}
\usepackage{float}
\usepackage{colortbl}
\usepackage{pdflscape}
\usepackage{tabu}
\usepackage{threeparttable}
\usepackage{threeparttablex}
\usepackage[normalem]{ulem}
\usepackage{makecell}
\usepackage{xcolor}
\usepackage{bookmark}
\IfFileExists{xurl.sty}{\usepackage{xurl}}{} % add URL line breaks if available
\urlstyle{same}
\hypersetup{
  pdftitle={EDA Check-in: Predicting Institutional Loan Distress},
  pdfauthor={MATH 456 Group Project},
  hidelinks,
  pdfcreator={LaTeX via pandoc}}

\title{EDA Check-in: Predicting Institutional Loan Distress}
\author{MATH 456 Group Project}
\date{April 17, 2026}

\begin{document}
\maketitle

\emph{Names:} Omar, Rowan, Christopher Duran

\subsection{Research question}\label{research-question}

Student loan defualts have become a heated topic among colleges and
students alike. This sparked the interest in researching predictive
features that could help classify or predict trends in colleges that
have many student loan defaults. Our response variable is the three-year
cohort default rate (\texttt{CDR3}), which is the fraction of an
institution's federal borrowers who defaulted within three years of
entering repayment. Our research question is: what institutional
features predict loan distress, after accounting for who the school
serves, how selective it is, what programs it offers, and how
well-resourced it is?

\subsection{The College Scorecard}\label{the-college-scorecard}

The College Scorecard is a public data set from the Department of
Education. Every year they publish one file per cohort (ex:
\texttt{MERGED2015\_16\_PP.csv}) with roughly 3,000 variables on every
post-secondary institution that takes federal financial aid. The numbers
come from three places stitched together: IPEDS for institutional
characteristics and enrollment, NSLDS for loan outcomes, and Treasury
tax records for post-graduation earnings. Because the data only covers
Title IV institutions and only describes students who received federal
aid, the population is specifically ``federally-aided students at
aid-participating schools'' --- not all college students nationally.

The raw data is massive and has way more variables than we can sensibly
use, so the first thing we did was cut it down to about 70 columns
organized into ten groups (identifiers, institutional characteristics,
admissions, student body, academics, cost, completion, debt, repayment,
earnings). From there we built a balanced panel by keeping only
institutions that showed up in every cohort year from 1996 through 2023
--- that gave us 3,790 schools. For this analysis we further filter to
rows where CDR3 is non-null, which narrows the window to 2011--2019.

Additionally, we cut 2020--2023 due to federal loan payments being
paused during COVID. Even though CDR3 technically exists for those
years, March 2020 through September 2023, by 2023 96\% of schools had
CDR3 pinned at exactly zero. Including those years could definitely lead
to less clarity and possible misleads.

After selecting our interval we come to 31,782 institution-year rows
across 3780 schools and 9 cohort years.

Further exploration lead to insight on net price (\texttt{NPT4}) being a
coalesce of the public and private net-price columns since they're
mutually exclusive by sector; same story for the 150\%-time completion
rate (\texttt{C150}). We log-transformed four right-skewed resource
variables (\texttt{UGDS}, \texttt{AVGFACSAL}, \texttt{INEXPFTE},
\texttt{TUITFTE}) and kept the raw columns too so we can pick whichever
works better downstream. And adjustments to some schools where made for
consistency, where we backfilled three time-invariant flags
(\texttt{HBCU}, \texttt{HSI}, \texttt{LOCALE}) due to schools
missingness in certain years even though they are constants.

\begingroup\fontsize{9}{11}\selectfont

\begin{longtable}[t]{>{\raggedright\arraybackslash}p{1.8cm}>{\raggedright\arraybackslash}p{4.5cm}>{\raggedright\arraybackslash}p{7cm}}
\caption{\label{tab:var-table}Variables we kept for the distress analysis.}\\
\toprule
Group & Variable & Description\\
\midrule
\cellcolor{gray!10}{Outcome} & \cellcolor{gray!10}{CDR3} & \cellcolor{gray!10}{3-yr cohort default rate}\\
Outcome & distressed & Binary: CDR3 > 0.10\\
\cellcolor{gray!10}{Identifier} & \cellcolor{gray!10}{UNITID, OPEID, INSTNM} & \cellcolor{gray!10}{Institution ID and name}\\
Structure & CONTROL & Public / Private NFP / Private FP\\
\cellcolor{gray!10}{Structure} & \cellcolor{gray!10}{PREDDEG} & \cellcolor{gray!10}{Predominant degree awarded}\\
\addlinespace
Structure & REGION, LOCALE & Census region and urbanicity\\
\cellcolor{gray!10}{Structure} & \cellcolor{gray!10}{HBCU, HSI, DISTANCEONLY} & \cellcolor{gray!10}{Institution-type flags}\\
Resources & INEXPFTE, AVGFACSAL, TUITFTE & Per-FTE spending / salary\\
\cellcolor{gray!10}{Resources} & \cellcolor{gray!10}{PFTFAC, STUFACR} & \cellcolor{gray!10}{Full-time faculty \%, student:faculty}\\
Admissions & ADM\_RATE, SAT\_AVG & Admission rate, median SAT\\
\addlinespace
\cellcolor{gray!10}{Students} & \cellcolor{gray!10}{UGDS, PCTPELL, PCTFLOAN} & \cellcolor{gray!10}{Enrollment, Pell / loan share}\\
Students & UGDS\_WHITE/BLACK/HISP/ASIAN & Race composition\\
\cellcolor{gray!10}{Students} & \cellcolor{gray!10}{FEMALE, UG25ABV} & \cellcolor{gray!10}{Female share, adult learner share}\\
Academics & PCIP01 ... PCIP54 & 38 CIP-2 program share columns\\
\cellcolor{gray!10}{Cost} & \cellcolor{gray!10}{NPT4} & \cellcolor{gray!10}{Avg net price (coalesced pub/priv)}\\
\addlinespace
Completion & C150 & 150\%-time completion rate\\
\cellcolor{gray!10}{Debt} & \cellcolor{gray!10}{GRAD\_DEBT\_MDN} & \cellcolor{gray!10}{Median graduate debt}\\
\bottomrule
\end{longtable}
\endgroup{}

\subsection{Observations}\label{observations}

Our approach had four rough phases. First we checked the distribution of
CDR3 itself across years to check consistency and/or trends. Then we
tried to figure out whether variation in CDR3 is mostly driven by
differences between schools or by changes within the same school over
time, since that tells us whether we need fancy panel methods or can
mostly get away with cross-sectional modeling. After that we looked at
how CDR3 relates to individual predictors one at a time, using boxplots
for the categorical ones and GAM smooths for the continuous ones.
Finally, we ran a correlation matrix across all the numeric predictors
to check for multicollinearity before starting the LASSO regression.

Everything below uses the tidyverse (\texttt{dplyr}, \texttt{tidyr},
\texttt{ggplot2}), with \texttt{lme4} and \texttt{performance} for the
mixed model that gives us the ICC.

\subsection{Results}\label{results}

\subsubsection{Descriptive statistics}\label{descriptive-statistics}

\begingroup\fontsize{9}{11}\selectfont

\begin{longtable}[t]{rrrrrr}
\caption{\label{tab:descriptives}Overall panel summary.}\\
\toprule
Observations & Schools & Mean CDR3 & Median CDR3 & SD CDR3 & \% Distressed (CDR3 > 0.10)\\
\midrule
\cellcolor{gray!10}{31782} & \cellcolor{gray!10}{3780} & \cellcolor{gray!10}{0.107} & \cellcolor{gray!10}{0.093} & \cellcolor{gray!10}{0.079} & \cellcolor{gray!10}{46.7}\\
\bottomrule
\end{longtable}
\endgroup{}

The average three-year default rate across the panel is around 10\%,
which is the historical ballpark. But the distribution is heavily
right-skewed: the median is lower than the mean, and a substantial chunk
of institution-years land above our 10\% ``distressed'' threshold. Table
2 breaks this down by sector and shows where the distress is really
concentrated.

\begingroup\fontsize{9}{11}\selectfont

\begin{longtable}[t]{lrrrrrr}
\caption{\label{tab:sector-table}CDR3 by institutional control sector.}\\
\toprule
Sector & Observations & Schools & Mean & Median & SD & \% Distressed\\
\midrule
\cellcolor{gray!10}{NA} & \cellcolor{gray!10}{31782} & \cellcolor{gray!10}{3780} & \cellcolor{gray!10}{0.107} & \cellcolor{gray!10}{0.093} & \cellcolor{gray!10}{0.079} & \cellcolor{gray!10}{46.7}\\
\bottomrule
\end{longtable}
\endgroup{}

\subsubsection{Time Distribution of
CDR3}\label{time-distribution-of-cdr3}

Figure 1 is just CDR3 faceted by year. The shape is right-skewed, modal
around 5--10\%, with a long upper tail. There's no big temporal shift in
the outcome across our nine-year window, which is good news because it
means year probably doesn't need to be treated as a big structural break
in the modeling. We can just use it as a regular predictor.

\begin{figure}[H]

{\centering \includegraphics[width=0.9\linewidth]{EDA_report_files/figure-latex/outcome-dist-1} 

}

\caption{Distribution of CDR3 by cohort year. The shape is pretty consistent across the nine-year window.}\label{fig:outcome-dist}
\end{figure}

\subsubsection{Between-School vs Within-School
Variance}\label{between-school-vs-within-school-variance}

This is the big structural question for our panel. If most of the
variance in CDR3 is between schools, meaning School A has always had a
certain default rate and School B has always had a different one, then
the problem is basically cross-sectional and we don't need to bother
with random effects. If there's meaningful within-school movement, then
the panel structure matters and we can exploit it with mixed models.

\begingroup\fontsize{9}{11}\selectfont

\begin{longtable}[t]{lrr}
\caption{\label{tab:variance-decomp}Manual variance decomposition of CDR3.}\\
\toprule
Component & Variance & \% of total\\
\midrule
\cellcolor{gray!10}{Between-school} & \cellcolor{gray!10}{0.00513} & \cellcolor{gray!10}{83}\\
Within-school & 0.00197 & 32\\
\cellcolor{gray!10}{Total} & \cellcolor{gray!10}{0.00618} & \cellcolor{gray!10}{100}\\
\bottomrule
\end{longtable}
\endgroup{}

The intraclass correlation from a random-intercept model confirms the
same story as the manual decomposition: ICC = 0.72. So about 72\% of the
total variance in CDR3 is stable between-school stuff, and only the
remaining 28\% is within-school year-to-year wobble. School identity is
the dominant factor, but there's enough within-school movement that a
mixed-model specification should still earn its keep.

Figure 2 shows what this looks like at the institution level. Each thin
line is one school's CDR3 over the nine cohort years. Schools pretty
clearly stay in their own lanes --- a low-default school today was a
low-default school eight years ago. The for-profit sector (right panel)
sits visibly higher than public and private non-profit institutions.

\begin{figure}[H]

{\centering \includegraphics[width=0.9\linewidth]{EDA_report_files/figure-latex/spaghetti-1} 

}

\caption{CDR3 trajectories 2011-2019 for a random sample of 75 schools per sector. Each line is one school.}\label{fig:spaghetti}
\end{figure}

Figure 3 plots each school's mean CDR3 against its within-school
standard deviation on a log-log scale. We switched to log-log because
the raw-scale version had everything clumped near the origin and was
impossible to read. The positive slope of the OLS fit tells us that
higher-default schools are also more volatile year-to-year, or a
mean-variance relationship, which is what you'd expect for a bounded
proportion like a default rate. With the for-profits cluster in the
upper right (high mean, high volatility), and publics in the lower left.

\begin{figure}[H]

{\centering \includegraphics[width=0.9\linewidth]{EDA_report_files/figure-latex/mean-vs-sd-1} 

}

\caption{School-level mean CDR3 vs within-school SD, log-log scale. Black line is OLS fit.}\label{fig:mean-vs-sd}
\end{figure}

\subsubsection{CDR3 \& Feature Relations}\label{cdr3-feature-relations}

Figure 5 is the categorical story. Public and private nonprofit
institutions sit at median CDR3 around 7--8\%, while for-profits are
visibly higher with a median near 15\% and an IQR that stretches well
above 20\%. Predominant degree matters a lot too --- certificate and
associate-focused institutions default more than bachelor's-granting
ones, which probably reflects the populations those programs serve and
the labor market outcomes of shorter credentials.

\begin{figure}[H]

{\centering \includegraphics[width=0.9\linewidth]{EDA_report_files/figure-latex/boxplots-1} 

}

\caption{CDR3 by institutional control sector and predominant degree awarded.}\label{fig:boxplots}
\end{figure}

Figure 6 shows CDR3 against five continuous predictors with cubic-spline
smooths through the cloud. A few things jump out. Pell Grant share has a
clean positive relationship with CDR3 --- schools serving more
low-income students default more, which is not surprising. Completion
rate has a strong negative relationship and it's pretty nonlinear, which
tells us a plain linear model probably isn't going to capture
everything. Median graduate debt is actually negatively associated with
default, which sounds backwards until you remember that the schools
where students borrow a lot are often selective private universities
serving wealthier students who then don't default. It's a selection
effect, not a causal claim. Net price shows a similar weak-negative
pattern for the same reason. These nonlinearities and counterintuitive
signs are a big part of why we want gradient boosting and GAMs in the
modeling stage instead of just OLS.

\begin{figure}[H]

{\centering \includegraphics[width=0.9\linewidth]{EDA_report_files/figure-latex/gam-plots-1} 

}

\caption{CDR3 vs five continuous predictors, with GAM cubic-spline smooths (red).}\label{fig:gam-plots}
\end{figure}

\subsubsection{Multi-Collinearity}\label{multi-collinearity}

Before throwing everything into a LASSO model, we wanted to see which
predictive features have a heavy correlation. Figure 7 is a Pearson
correlation matrix across all the numeric predictors, reordered by
hierarchical clustering so the correlated blocks sit next to each other.
Most of the off-diagonal cells are near zero, which was honestly a bit
surprising --- we expected more redundancy. The different dimensions
we're measuring (selectivity, cost, demographics, program mix, faculty
resources) mostly seem to capture genuinely different things about a
school rather than being proxies for each other.

\begin{figure}[H]

{\centering \includegraphics[width=0.9\linewidth]{EDA_report_files/figure-latex/corr-heatmap-1} 

}

\caption{Pearson correlation across continuous predictors, reordered by hierarchical clustering.}\label{fig:corr-heatmap}
\end{figure}

Only a handful of pairs exceed our \(|r| > 0.8\) flag. Those are the
ones we'll need to consolidate or drop before LASSO, since penalized
regression gets unstable when it has to pick between near-duplicate
predictors. But, a short list means we don't have to do aggressive
pruning.

\begingroup\fontsize{9}{11}\selectfont

\begin{longtable}[t]{llr}
\caption{\label{tab:high-corr-table}Predictor pairs with |r| > 0.8, flagged for LASSO pruning.}\\
\toprule
var1 & var2 & r\\
\midrule
\cellcolor{gray!10}{AVGFACSAL} & \cellcolor{gray!10}{AVGFACSAL\_log} & \cellcolor{gray!10}{0.958}\\
\bottomrule
\end{longtable}
\endgroup{}

\subsection{What's Next?}\label{whats-next}

The EDA leaves us with a reasonably clear setup for the modeling phase.
CDR3 is well-behaved and we have a clean nine-year window. School
identity is the biggest single driver of variation, but there's enough
within-school movement that a hierarchical specification still adds
something. Several of the bivariate relationships are nonlinear or
counterintuitive, which tells us we want flexible models like XGBoost
alongside the linear baselines. And the feature space is mostly
orthogonal so LASSO selection should be stable.

Our plan from here is a model ladder: a null baseline, then a small OLS
with hand-picked predictors for interpretability, then penalized
regression (LASSO and ElasticNet), then XGBoost with SHAP-based
interpretation for the nonlinearities, and finally a linear mixed model
to quantify how much of the explanatory power really comes from school
identity versus measured institutional features. We'll also run the
whole thing with the binary \texttt{distressed} outcome as a robustness
check and use \texttt{RPY\_3YR\_RT} on its overlap window as a secondary
outcome.

\(\textbf{Note:}\) A few caveats worth naming up front. Everything in
the data is institution-level, so we can't say anything about individual
students. The outcome is only defined for students who received federal
aid, so schools with low Title IV participation are underrepresented.
And the privacy suppression on small cells means very small institutions
are systematically missing --- probably correlated with the variables we
care about most. These aren't things we can fix in the modeling; they
just shape how far we can push the conclusions.

\end{document}
